% Преамбула сборника. Подключается из main.tex и из сборок отдельных глав.
\usepackage[T2A]{fontenc}     % кодировка шрифтов
\usepackage[utf8]{inputenc}   % кодировка исходного файла
\usepackage[russian]{babel}   % поддержка русского
\usepackage[a4paper,margin=1.2cm]{geometry}
\usepackage{amsthm, amssymb}
\usepackage{amsmath}
\usepackage{float}
\usepackage{graphicx} % Required for inserting images
% Иллюстрации разложены по происхождению, но в \includegraphics по-прежнему
% указывается только имя файла.
\graphicspath{{figures/book/}{figures/own/}}
\usepackage{caption}
\usepackage{wrapfig}
\usepackage{subcaption}
\usepackage{tikz}
\usetikzlibrary{positioning,arrows.meta}
% --- навигация по документу -------------------------------------------------
% hyperref загружается до algorithm2e: обратный порядок ломает нумерацию
% строк в окружении algorithm.
\usepackage[unicode, hidelinks]{hyperref}
\hypersetup{
    bookmarksnumbered = true,
    bookmarksopen     = true,
    pdftitle          = {Ответы к книге «Обучение с подкреплением», Саттон и Барто},
    pdfsubject        = {Решения упражнений},
}

\usepackage[ruled, linesnumbered]{algorithm2e}

\captionsetup{labelformat=empty,hypcap=false}
\emergencystretch=3em
\hfuzz=30pt
\hbadness=10000
\sloppy

\renewcommand{\qedsymbol}{$\blacksquare$}

\newenvironment{solution}
  {\begin{proof}[Решение]}
  {\end{proof}}
